Protein language models encode residue-level regularities that support structure, function, and mutation-effect prediction, but it remains unclear when internal features should be treated as new biological hypotheses rather than rediscoveries of known annotations or artifacts of high-dimensional search. We examine this question in a compact, reproducible workflow that combines sparse feature association, motif exclusion, baseline controls, and internal intervention. We run \esmeight with pretrained layer-4 \interplm sparse autoencoder features on 16 short \proteingym deep mutational scanning assays, covering 12,085 single substitutions and 688 residue positions. The strongest sparse feature per assay correlates with residue-level mutation sensitivity, with median absolute Spearman correlation 0.547. However, this does not significantly exceed raw \esm hidden dimensions (median 0.544; one-sided Wilcoxon $p=0.798$) or shuffled sparse-feature controls (median shuffled 95th percentile 0.543; $p=0.217$). We identify three low \elm/\prosite alignment feature-assay candidates. One exploratory intervention, ablating \interplm feature 6419 in \texttt{YNZC\_BACSU\_Tsuboyama\_2023\_2JVD}, changes \esm masked-marginal variant scores more at high-activation positions than at low-activation positions (effect 0.0223 log-prob units; Mann-Whitney $p=0.00495$). These results do not establish a novel biological mechanism. They instead support a cautious triage protocol for turning protein-model features into testable hypotheses.
